\section{Osnove kompleksne analize} \label{sec:hipgeom}

V tem poglavju ponovimo nekaj osnovnih oznak in definicij iz kompleksne analize, ki jih bomo uporabljali kasneje. Nato obravnavamo kompleksni eksponent in logaritem, za tem pa uvedemo še \emph{hiperbolično metriko} in preučimo njene lastnosti.

\subsection{Osnove holomorfnosti}

Kompleksno število \(z \in \CC\) lahko zapišemo kot vsoto \(z = x + i y\) ali v polarnem zapisu \(z = r e^{i \theta} = r (\cos \theta + i \sin \theta)\). Kot, ki ga oklepata abscisa in premica skozi \(0\) in \(z\), označimo z \(\Arg (z) \in (- \pi, \pi]\) ali \(\arg (z) \in [0, 2 \pi)\). Z \(\Delta (a, r)\) označimo odprt disk s središčem v \(a\) in radijem \(r\). Za enotski disk nadalje uporabimo oznako \(\DD \coloneq \Delta (0, 1)\).

Neprazno, odprto in povezano podmnožico kompleksnih števil \(D \subseteq \CC\) imenujemo \emph{območje}. Večinoma se bomo ukvarjali s kompleksnimi funkcijami ene spremenljivke \(f \colon D \to \CC\). Le-te lahko zapišemo kot vsoto realnega in imaginarnega dela \(f = u + iv\), kjer sta \(u = u (x, y)\) in \(v = v (x, y)\) ustrezni realni funkciji. Naj bo od tu naprej \(U \subseteq \CC\) odprta množica.

\begin{definicija}
    Pravimo, da je funkcija \(f \colon U \to \CC\) \emph{holomorfna} v \(z \in U\), če obstaja limita
    \[\lim_{w \to z} \frac{f (w) - f (z)}{w - z} = \lim_{\zeta \to 0} \frac{f (z + \zeta) - f (z)}{\zeta}.\]
\end{definicija}

Holomorfnost razumemo kot odvedljivost v kompleksnem smislu. Zato zgornji limiti pravimo \emph{odvod funkcije \(f\) v točki \(z\)} in jo označimo s \(f' (z)\). Kot za realne veljajo tudi za kompleksne funkcije običajna pravila za odvajanje. Poleg tega pa nam holomorfnost karakterizira tudi naslednji izrek. Dokaz najdemo v knjigi \cite[izrek 43]{Globevnik_2010}.

\begin{izrek} \label{thm:cr}
    Funkcija \(f = u + i v\) je na množici \(U\) holomorfna natanko tedaj, ko sta \(u\) in \(v\) na množici \(\set{(x, y) \in \RR^2 : x + i y \in U}\) zvezno parcialno odvedljivi in zadoščata \emph{Cauchy-Riemannovemu sistemu enačb}
    \begin{equation} \tag{CR} \label{eq:cr}
        \frac{\partial u}{\partial x} = \frac{\partial v}{\partial y}, \qquad \frac{\partial u}{\partial y} = - \frac{\partial v}{\partial x}.
    \end{equation}
\end{izrek}

% \begin{dokaz}
%     Naj bo \(f = u + i v\), \(z_0 = x + i y\) in \(\zeta = h + i k\). Od \(f\) lahko, ne da bi to vplivalo na odvedljivost, odštejemo konstanto \(f (z_0)\). Zato privzamemo, da velja \(f (z_0) = 0\). Če krajše pišemo \(u = u (x + h, y + k) - u (x, y)\) in \(v = v (x + h, y + k) - v (x, y)\), velja
%     \[\frac{f (z_0 + \zeta)}{\zeta} = \frac{u + i v}{h + i k} = \frac{u h + v k}{h^2 + k^2} + i \frac{v h - u k}{h^2 + k^2}.\]
%     Uvedemo polarne koordinate \(h = i \cos \varphi\) in \(k = r \sin \varphi\), v katerih velja
%     \[\frac{h^2}{h^2 + k^2} = \cos^2 \varphi, \quad \frac{k^2}{h^2 + k^2} = \sin^2 \varphi.\]
%     S tem se zgornji izraz poenostavi v
%     \[\frac{f (z_0 + \zeta)}{\zeta} = \frac{u}{h} \cos^2 \varphi + \frac{v}{k} \sin^2 \varphi + i \prt{\frac{v}{h} \cos^2 \varphi - \frac{u}{k} \sin^2 \varphi}.\]
%     V izraz na ustreznih mestih vrinemo parcialne odvode ter poenostavimo, da dobimo
%     \[
%         \begin{multlined}[15cm]
%             \frac{f (z_0 + \zeta)}{\zeta} = \prt{\frac{u}{h} - u_x} \cos^2 \varphi + \prt{\frac{v}{k} - v_y} \sin^2 \varphi + u_x + (v_y - u_x) \sin^2 \varphi\\
%             + i \left[ \prt{\frac{v}{h} - v_x} \cos^2 \varphi - \prt{\frac{u}{k} - u_y} \sin^2 \varphi + v_x - (u_y + v_x) \sin^2 \varphi \right].
%         \end{multlined}
%     \]
%     Ko gre \(\zeta \to 0\), gredo po definiciji parcialnega odvoda proti \(0\) tudi razlike v oklepajih. Vidimo, da bo limita \(\zeta \to 0\) izraza na levi obstajala natanko tedaj, ko bo veljalo \(v_y - u_x = u_y + v_x = 0\). V tem primeru bo \(f' (z_0) = u_x + i v_x\).
% \end{dokaz}

Za holomorfne funkcije, ki fiksirajo izhodišče, velja naslednja lema. Njen dokaz najdemo v knjigi \cite[lema 1]{Globevnik_2010}.

\begin{lema}[Schwarz] \label{lem:schwarz}
    Naj bo \(f \colon \DD \to \DD\) holomorfna funkcija, za katero velja \(f (0) = 0\). Potem drži natanko ena od naslednjih trditev.
    \begin{enumerate}
        \item Velja \(|f (z)| < |z|\) za \(z \neq 0\) in \(|f' (0)| < 1\).
        \item Obstaja \(\theta \in \RR\), da za vsak \(z\) velja \(f (z) = e^{i \theta}\) in \(|f' (0)| = |e^{i \theta}| = 1\).
    \end{enumerate}
\end{lema}

Če je \(f\) holomorfna na \(\CC\), pravimo, da je \emph{cela funkcija}. Vemo, da je to v \(\CC\) ekvivalentno pogoju, da je \(f\) kompleksno analitična pri vsakem \(z \in \CC\). Cela funkcija je \emph{transcendentna}, če ni polinom ali kvocient dveh polinomov. Pravimo, da je funkcija \emph{konformna}, če ohranja kote in smisel vrtenja. Vemo, da so holomorfne funkcije konformne natanko tedaj, ko je njihov odvod neničelen (glej \cite[izrek 78]{Globevnik_2010}).

Funkcija \(f \colon U \to U'\) je \emph{biholomorfna}, če je bijektivna, holomorfna in ima holomorfen inverz \(f^{-1} \colon U' \to U\). Če je \(D\) območje in \(f \colon D \to D'\) bijektivna in holomorfna, vemo, da je biholomorfna (glej \cite[izrek 77]{Globevnik_2010}). Območje je \emph{enostavno povezano}, če je homeomorfno enotskemu disku. Topološko definicijo nam s kompleksno lastnostjo poveže naslednji znamenit izrek.

\begin{definicija}[Riemannov upodobitveni izrek]
    Vsako enostavno povezano območje različno od \(\CC\) je biholomorfno enotskemu disku.
\end{definicija}

\subsection{Kompleksni eksponent in logaritem}

\noindent Za \(z = x + i y \in \CC\) je kompleksna eksponentna preslikava podana s predpisom
\[e^z = e^x (\cos y + i \sin y).\]
Označimo \(u (x, y) = e^x \cos y\) in \(v (x, y) = e^x \sin y\) in izračunamo
\[\frac{\partial u}{\partial x} (x, y) = e^x \cos y = \frac{\partial v}{\partial y} (x, y), \qquad \frac{\partial u}{\partial y} (x, y) = - e^x \sin y = -\frac{\partial v}{\partial x} (x, y).\]
Po izreku \ref{thm:cr} je kompleksni eksponent holomorfen na \(\CC\), torej cela funkcija. Dodatno pa je tudi konformna funkcija, saj je odvod neničelen. Ker je
\[e^{z + 2 \pi i} = e^x (\cos (y + 2 \pi) + i \sin (y + 2 \pi)) = e^z,\]
je tudi \(2\pi i\)-periodičen in posledično nima inverza. Res, naj bo \(w = u + i v\). Rešimo enačbo:
\begin{align*}
    e^w = z & \iff e^u (\cos v + i \sin v) = x + i y                                    \\
            & \iff e^u = \sqrt{x^2 + y^2} \quad \text{in} \quad v = \Arg (z) + 2 k \pi.
\end{align*}
Torej je \(e^w = z\) natanko tedaj, ko je
\[w = \ln |z| + i \Arg (z) + 2 k \pi i\]
za neki \(k \in \ZZ\). Prav tako bi lahko namesto \(\Arg\) izbrali \(\arg\). Da lahko definiramo inverz, moramo domeno omejiti. Naj bo \(\theta \in \RR\). Potem je
\[\begin{multlined}[10cm]
        \exp \colon  \set{\zeta \in \CC \colon \theta -\pi < \im \zeta < \theta + \pi}\\
        \to \set{\omega \in \CC \colon \Arg (\omega) \not\equiv \theta - \pi\ (\operatorname{mod}\, 2 \pi)} \eqcolon U
    \end{multlined}\]
bijektivna na \(U\) in ima tam holomorfen inverz \(L \colon U \to \CC\):
\[L (\omega) \coloneq \log |\omega| + i \cdot \prt{\theta + \Arg \prt{\frac{\omega}{e^{i \theta}}}}.\]
Tako je kompleksni eksponent na zgoraj definiranih območjih biholomorfen. Vsaki taki funkciji \(L\) rečemo \emph{veja logaritma}. Če vzamemo \(\theta = 0\) dobimo \emph{glavno} vejo logaritma. Veje logaritma obstajajo tudi na vsakem enostavno povezanem območju, vendar bodo za nas pomembne le zarezane ravnine \(\CC \setminus [0, \infty)\) in diski, ki ne vsebuje izhodišča. Na primer, imejmo disk \(\Delta \subset \CC\), tako da \(0 \notin \Delta\). Naj bo \(\zeta_0 \in \CC\), da je \(e^{\zeta_0} \in \Delta\). Če vzamemo \(\theta = \im \zeta_0\), dobimo holomorfno preslikavo \(L \colon \Delta \to \CC\), za katero velja \(L (e^{\zeta_0}) = \zeta_0\) ter \(e^{L (\omega)} = \omega\) in \(\im \zeta_0 - \pi < \im L (\omega) < \im \zeta_0 + \pi\) za vsak \(\omega \in \Delta\).

\begin{figure}
    \centering
    \includegraphics[width=1\textwidth]{needham_figure.pdf}
    \caption[Eksponentne preslikave]{Eksponentna preslikava, po knjigi \cite{Needham_1997}.}
    \label{fig:exponential}
\end{figure}
Dobro vizualizacijo kompleksnega eksponenta nam nudi slika \ref{fig:exponential}, na kateri je \(w = e^z\). Na neformalen način navedemo naslednja opažanja.
\begin{itemize}
    \item Če točka \(z\) potuje navzgor s hitrostjo \(v\), potem se točka \(w\) vrti okrog izhodišča s kotno hitrostjo \(v\). Ko \(z\) prepotuje razdaljo \(2 \pi\), \(w\) enkrat zaokroži okoli izhodišča. Torej je preslikava periodična s periodo \(2 \pi i\).
    \item Če točka \(z\) potuje proti levi s konstantno hitrostjo, točka \(w\) potuje proti izhodišču s pojemajočo hitrostjo. Obratno, če \(z\) potuje proti desni, potem \(w\) potuje stan od izhodišča s pospešeno hitrostjo.
    \item Po prejšnjih dveh alinejah bo celo \(w\)-ravnino brez točke \(w = 0\) pokril vodoraven pas višine \(2 \pi\) v \(z\)-ravnini.
    \item Vodoravne premice se preslikajo v žarke iz izhodišča, v posebnem se realna os preslika v \([0, \infty)\). Navpične črte se preslikajo v krožnice okrog izhodišča, v posebnem se imaginarna os preslika v enotsko krožnico. Splošna premica se preslika v spiralo, ki se navija okrog izhodišča.
    \item Sivo obarvana leva polravnina se preslika v enotski disk, desna polravnina pa v njegov komplement.
    \item Slike majhnih koordinatnih kvadratov ostanejo homeomorfne kvadratom, pri čemer se krožnice s poltraki sekajo pod pravim kotom. To ponazarja, da je preslikava konformna---ohranja kote.
\end{itemize}

\subsection{Hiperbolična geometrija}

Temelj evklidske geometrije nam postavi pet Evklidovih aksiomov.
\begin{enumerate}
    \item Skozi poljubni dve točki poteka natanko ena premica.
    \item Premica je neomejena.
    \item Za vsako daljico obstaja krožnica, ki ima to daljico za polmer in eno od krajišč za središče.
    \item Vsi pravi koti so skladni.
    \item Za poljubno premico \(p\) in točko \(A\), ki ne leži na premici \(p\), obstaja natanko ena premica \(q\), ki vsebuje \(A\) in ne seka premice \(p\).
\end{enumerate}

\noindent Zadnji peti aksiom se imenuje \emph{aksiom o vzporednici}. V hiperbolični geometriji ga nadomestimo z

\begin{enumerate}[resume]
    \item[5'.] Obstajata premica \(p\) in točka \(A\), ki ne leži na premici \(p\), tako da obstajata vsaj dve premici \(q\) in \(r\), ki vsebujeta \(A\) in ne sekata premice \(p\).
\end{enumerate}

\noindent Nov aksiom prinaša veliko posledic, med drugim tudi drugo metriko; \emph{hiperbolično metriko}, katere lastnosti bomo uporabili pri dokazu. Metriko najprej definiramo na enotskem disku \(\DD\) in jo nato razširimo na bolj splošne množice. Poglavje je povzeto po članku \cite[razdelek 3]{Shen_2015}.

Hiperbolično metriko podamo prek \emph{hiperboličnega ločnega elementa} najprej na enotskem disku:
\[\dd \rho_{\DD} (z) = \frac{2 |\dd z|}{1 - |z|^2}.\]
Hiperbolični ločni element nam pove, da se infinitezimalna sprememba v točki \(z\) v hiperbolični metriki izrazi kot infinitezimalna sprememba v evklidski metriki, pomnožena s t.i.~\emph{gostoto} hiperbolične metrike
\[\rho_{\DD} (z) = \frac{2}{1 - |z|^2}.\]
Iz tega izhaja naslednja definicija.

\begin{definicija}
    Naj bo \(\gamma \colon [a, b] \to \DD\) odsekoma \(C^1\) krivulja. Njena \emph{hiperbolična dolžina} je
    \[l_{\DD} (\gamma) \coloneq \int_{\gamma} \rho_{\DD} (z) \, | \dd z | = \int_{a}^{b} \frac{2 | \gamma' (t) | }{1 - | \gamma (t) |^2} \, \dd t.\]
\end{definicija}

\noindent Končno hiperbolično metriko definiramo kot sledi.

\begin{definicija}
    \emph{Hiperbolična metrika} je preslikava \(d_{\DD} \colon \DD \times \DD \to [0, \infty)\), definirana kot
    \[d_{\DD} (z, w) \coloneq \inf_{\gamma} l_{\DD} (\gamma),\]
    kjer \(\gamma\) teče po vseh odsekoma \(C^1\) krivuljah v \(\DD\), ki povezujejo točki \(z\) in \(w\).
\end{definicija}

Definicijo nam iz enotskega diska na enostavno povezana območja prenese naslednji izrek, katerega dokaz najdemo v delu \cite[izrek 6.4]{Beardon_2006}.

\begin{izrek}[Pick] \label{thm:pick}
    Naj bosta \(U, V \subsetneq \CC\) enostavno povezani območji in \(f \colon U \to V\) holomorfna preslikava. Potem na \(U\) obstaja enolično določena gostota hiperbolične metrike \(\rho_U\), da veljajo naslednje izjave.
    \begin{enumerate}
        \item Za vsak \(z \in \DD\) velja \(\rho_{\DD} (z) = \frac{2}{1 - |z|^2}\).
        \item Preslikava \(f\) ne veča \emph{hiperboličnega odvoda}, to pomeni, da velja \[\hder{f (z)}{U}{V} \coloneq |f' (z)| \cdot \frac{\rho_V (f (z))}{\rho_U (z)} \leq 1.\]
        \item Enakost \(\hder{f (z)}{U}{V} = 1\) velja natanko tedaj, ko je \(f\) biholomorfna.
        \item Če je \(U \subsetneq V\), potem za vsak \(z \in U\) velja \(\rho_U (z) > \rho_V (z)\).
    \end{enumerate}
\end{izrek}

\noindent S pomočjo druge in tretje točke lahko izračunamo gostoto hiperbolične metrike eksplicitno. To je prikazano v naslednji trditvi.

\begin{trditev}[Zgledi hiperboličnih metrik] \label{prop:hypexamples} \mbox{}
    \begin{enumerate}
        \item Za desno polravnino \(\HH \coloneq \set{ z \in \CC : \re (z) > 0 }\) je \(\displaystyle \rho_\HH = \frac{1}{\re (z)}\).
        \item Za zgornjo polravnino \(\KK \coloneq \set{z \in \CC : \im (z) > 0}\) je \(\displaystyle \rho_{\KK} (z) = \frac{1}{\im (z)}\)
        \item Za pas \(S \coloneq \set{z \in \CC : |\im (z)| < \pi}\) je \(\displaystyle \rho_{S} = \frac{1}{2 \cos (\im (z) / 2)}\).
        \item Za zarezani ravnini velja
              \[
                  \rho_{\CC \setminus [0, \infty)} (z) = \frac{1}{2 |z| \sin (\arg (z) / 2)},
                  \qquad
                  \rho_{\CC \setminus (-\infty, 0]} (z) = \frac{1}{2 |z| \cos (\Arg (z) / 2)}.
              \]
    \end{enumerate}
\end{trditev}

\begin{dokaz}
    Uporabimo izrek \ref{thm:pick} na štirih biholomorfnih funkcijah. Naj bo \(\varphi_1 \colon \HH \to \DD\); \(z \mapsto (1 - z) / (1 + z)\). Potem je
    \[\rho_{\HH} (z) = \rho_{\DD} (\varphi_1 (z)) \cdot \abs{\varphi_1' (z)} = \frac{2}{1 - |\varphi_1 (z)|^2} \cdot \frac{2}{|z + 1|^2} = \frac{4}{|1 + z|^2 - |1- z|^2}.\]
    Za \(z = x + i y\) razpišemo \(|1 \pm z|^2 = 1 \pm 2x + x^2 + y^2\) in dobimo želen rezultat.

    Naj bo \(\varphi_2 (z) \colon \KK \to \DD\); \(\varphi_2 (z) = (z - i) / (z + i)\). Potem po zgornjem postopku dobimo
    \[\rho_{\KK} (z) = \rho_{\DD} (\varphi_2 (z)) \cdot \abs{\varphi_2' (z)} = \frac{4}{|z + i|^2 - |z - i|^2} = \frac{1}{\im (z)}.\]

    \noindent Naj bo \(\varphi_3 \colon S \to \HH\); \(z \mapsto e^{z / 2}\). Potem je
    \[\rho_S (z) = \rho_{\HH} (\varphi_3 (z)) \cdot |\varphi_3' (z)| = \frac{|\varphi_3 (z)|}{2 \re \varphi_3 (z)} = \frac{1}{2 \cos (\arg (\varphi_3 (z)))} = \frac{1}{2 \cos (\im (z) / 2)}.\]
    Zadnja enakost drži, saj iz \(|\im (z)| < \pi\) sledi \(\arg (\varphi_3 (z)) = \im (z) / 2\).

    Naj bo \(\varphi_4 \colon \CC \setminus (- \infty, 0] \to \HH\); \(z \mapsto \sqrt{z}\) glavna veja korena. Ta funkcija je biholomorfna, če se omejimo na rešitev \(r e^{i \theta} \mapsto \sqrt{r} e^{i \theta / 2}\) za \(\theta \in (- \pi, \pi)\). Potem je
    \[\rho_{\CC \setminus (- \infty, 0]} (z) = \rho_{\HH} (\varphi_4 (z)) \cdot |\varphi_4' (z)| = \frac{1}{\re (\sqrt{z})} \cdot \frac{1}{2 |\sqrt{z}|}.\]
    Dobljeno poenostavimo s polarnim zapisom, da dobimo
    \[\frac{1}{2 r \cos (\theta / 2)} = \frac{1}{2 |z| \cos (\Arg (z) / 2)}.\]
    % \[\rho_{\CC \setminus [0, \infty)} (z) = \rho_{\HH} (\varphi_2 (z)) \cdot |\varphi_2' (z)| = \frac{2 |z|}{\re (- z^2)}.\]
    % Uporabimo polarni zapis \(z = r e^{i \theta}\), da dobimo \(- z^2 = - r^2 e^{2 i \theta}\) in \(\re (- r^2 e^{2 i \theta}) = - r^2 \cos (2 \theta)\). Zato z uporabo formule \(\cos (2 \theta) = \cos^2 \theta - \sin^2 \theta\) dobimo
    % \[\frac{2 |z|}{\re (- z^2)} = \frac{2}{- r \cos (2 \theta)} = \frac{2}{r (\sin^2 \theta - \cos^2 \theta)}\]
\end{dokaz}
